\subsection{Receiver operation and input boundaries}
\label{sec:receiver-operation}

The receiver turns the observed reply locations and comment topics back into the sender's message. Its work has two parts: establishing the same interpretation of the public context, then extracting the choices from the carrier comments. This section brings those operations together; Section~\ref{sec:multiframe} defines the multi-comment transport fields.

\subsubsection{Public observations, shared setup, and private data}
The observable inputs are the source post, its comments, and the carrier text with its comment and parent identifiers. The parties also agree on the preparation rules: models and prompts, source selection, comment cutoff, list ordering, compression representation, and transport and protection settings. A keyed protection mode additionally requires the shared secret. These agreements determine how the same visible choice maps to bits.

The topic list and compression dictionary are different objects. The first supplies the alternatives used to interpret a comment's topic; the second interprets compressed text references after the bits have been recovered. Both must match the sender's corresponding structures. The receiver can construct them through the agreed preparation process, or receive prepared structures when the experiment specifically uses that condition. The availability of prepared inputs is stated separately from the decoding algorithm.

The original message, intended topic indices, and sender acceptance records are private evaluation data. Independent extraction infers the choices from the observations rather than taking them from those records. The historical audit-assisted comparison is a separate execution path, described in Section~\ref{sec:evaluation}. For multi-comment messages, the production receiver accepts an ordered list of carrier references. In the archive experiments, a separate driver constructs these references from public timestamps and identifiers under the agreed timing rule.

\subsubsection{From the observed context to the recovered message}
\begin{enumerate}
\item \textbf{Prepare the receiving context.} The reconstruction path performs data loading, domain research, and topic generation using the agreed inputs and settings. Carrier comments are excluded from the pre-embedding context. The shared codec constructs the compression dictionary. Identity hashes support comparison of structures during validation; they do not replace the structures needed for decoding.
\item \textbf{Read each carrier's two choices.} The parent identifier locates the reply target in the ordered list. For the topic choice, the decoder retrieves nearby topics with semantic search, applies lexical reranking, and asks the identification model to select an index. It uses the same list and choice-mapping rules as the sender. The single-comment and multi-comment paths retain their different width rules described above.
\item \textbf{Assemble and check the bits.} In multi-comment mode, the receiver joins candidate slices in carrier order and reads the marker, version, count, and length. It checks the payload CRC and final padding and requires one recoverable payload. A missing carrier, unresolved interpretation, or failed integrity check produces a failure result rather than a partial message reported as complete.
\item \textbf{Recover the message representation.} The receiver reverses the agreed payload transform and decompresses as required by that representation. Its output includes the recovered text and a success or failure status, with diagnostic information identifying the stage reached. Evaluation compares the recovered bytes with the private original only after extraction.
\end{enumerate}

\subsubsection{Component validation and complete-message experiments}
Receiver validation is supported by several kinds of evidence. Unit and component tests exercise parent lookup, removal of carrier comments from context, reconstruction orchestration, choice interpretation, ordered stream assembly, exact compression round trips, and rejection of invalid references or corrupted streams. Controlled substitutes for external acquisition and model responses isolate these operations, so a test can establish whether the receiver handles a given observation correctly without another network or model call.

The message experiments then exercise combined operations with generated carrier text. The source-excerpt experiment tests topic identification and complete-message extraction with prepared shared inputs. The archive experiments add independent public-context acquisition and list reconstruction for the deterministic thread-only variant. Development checks separately compare repeated preparation of search queries, documents, and topics (Section~\ref{sec:reconstruction-validation}). Thus the receiver has component validation and complete-message experimental evidence. The remaining integration scope is a single end-to-end experiment combining full web/model-based preparation with actual platform publication and reception; that scope is distinct from whether the individual receiver components have been tested.
